\noindent\textit{Phạm vi.} Lời giải cho 22 bài được chỉ định trong Johnson--Wichern,
\textit{Applied Multivariate Statistical Analysis}, ấn bản 6, Chương 4--5.
Các nhãn bài theo sách; phần đề chỉ ghi tóm tắt, các phép tính và chứng minh
được trình bày độc lập. Dữ kiện được đối chiếu với bản sách, không suy ra
chỉ từ sự tương đồng với Exercise 1--2.

\medskip
\noindent\textbf{Quy ước và công thức dùng chung.}
Với mẫu $\mathbf{x}_1,\ldots,\mathbf{x}_n\in\mathbb R^p$, đặt
\[
\bar{\mathbf{x}}=\frac1n\sum_i\mathbf{x}_i,\qquad
\mathbf S=\frac1{n-1}\sum_i(\mathbf{x}_i-\bar{\mathbf{x}})
(\mathbf{x}_i-\bar{\mathbf{x}})'.
\]
Các phân vị $F_{a,b}(q)$, $t_\nu(q)$ và $\chi_p^2(q)$ là phân vị tích lũy mức $q$.
Với mẫu chuẩn, $\boldsymbol\Sigma>0$ và $n>p$, dưới $H_0:\boldsymbol\mu=\boldsymbol\mu_0$,
\[
T^2=n(\bar{\mathbf x}-\boldsymbol\mu_0)'\mathbf S^{-1}
(\bar{\mathbf x}-\boldsymbol\mu_0),\qquad
\frac{n-p}{p(n-1)}T^2\sim F_{p,n-p}.
\]
Đặt $c_T=\frac{p(n-1)}{n-p}F_{p,n-p}(1-\alpha)$. Miền tin cậy là
$n(\bar{\mathbf x}-\boldsymbol\mu)'\mathbf S^{-1}(\bar{\mathbf x}-\boldsymbol\mu)\le c_T$.
Khoảng đồng thời cho mọi tổ hợp tuyến tính là
\[
\mathbf a'\bar{\mathbf x}\pm\sqrt{c_T\mathbf a'\mathbf S\mathbf a/n}.
\]
Với $m$ đại lượng được chỉ định trước, Bonferroni dùng
$t_{n-1}(1-\alpha/(2m))$ thay cho $\sqrt{c_T}$.
Khi chỉ có giả thiết mẫu lớn và số chiều cố định, dùng $\chi_p^2(1-\alpha)$
thay cho $c_T$; tính chính xác của phân phối $F$ đòi hỏi mẫu chuẩn.
